%~Mouliné par MaN_auto v.0.23.0 2020-05-19 15:13:25
\documentclass[AHL,Unicode,longabstracts]{cedram}

\makeatletter
\def\@altabstractlanguage{}
\makeatother


\AtBeginDocument{
\newtheorem{fact}[cdrthm]{Fact} }



\newtheorem{thmx}{Theorem}
\newtheorem{corox}[thmx]{Corollary}
\renewcommand*{\thethmx}{\Alph{thmx}}

%%%%%%%%%%Mathop%%%%%%%%%%%
\DeclareMathOperator{\ad}{ad}
\DeclareMathOperator{\Ad}{Ad}
\DeclareMathOperator{\Aut}{Aut}
\DeclareMathOperator{\Conf}{Conf}
\DeclareMathOperator{\GL}{GL}
\DeclareMathOperator{\Hom}{Hom}
\DeclareMathOperator{\iso}{Iso}
\DeclareMathOperator{\Iso}{Iso}
\DeclareMathOperator{\Ker}{Ker}
\DeclareMathOperator{\Nor}{Nor}
\DeclareMathOperator{\PSL}{PSL}
\DeclareMathOperator{\PGL}{PGL}
\DeclareMathOperator{\SL}{SL}
\DeclareMathOperator{\Sol}{SOL}
\DeclareMathOperator{\SO}{SO}
\DeclareMathOperator{\ein}{{\bf Ein}_3}

\DeclareMathOperator{\Omp}{O}%%\Omp
\DeclareMathOperator{\sol}{sol}

\DeclareMathOperator{\Span}{Span}
\DeclareMathOperator{\Sop}{S}
\DeclareMathOperator{\vol}{vol}
\DeclareMathOperator{\rk}{rk}
\DeclareMathOperator{\flatm}{flat}
\DeclareMathOperator{\Is}{Is}
\DeclareMathOperator{\Homeo}{Homeo}
\DeclareMathOperator{\hol}{hol}
\DeclareMathOperator{\Diff}{Diff}
\DeclareMathOperator{\Trace}{Trace}
\DeclareMathOperator{\Fix}{Fix}
\DeclareMathOperator{\rank}{rank}
\DeclareMathOperator{\loc}{loc}
\DeclareMathOperator{\id}{id}
\DeclareMathOperator{\Id}{Id}





\newcommand{\mint}{{M^{\intl}}}
\newcommand{\hmint}{{{\hat M}^{\intl}}}
\DeclareMathOperator{\intl}{int}


\DeclareMathOperator{\slop}{sl}
\newcommand{\sld}{\slop(2,\RR)}
\DeclareMathOperator{\kil}{{\mathfrak{kill}}}
\newcommand{\kiloc}{\kil^{\loc}}
\DeclareMathOperator{\heis}{{\mathfrak{heis}}}
\DeclareMathOperator{\Heis}{Heis}
\DeclareMathOperator{\SOL}{SOL}


%%%%%%%%%%%%%%%%%%%%%%%MATHCAL%%%%%%%%%%%%
\newcommand{\cald}{\mathcal{D}}
\newcommand{\calf}{\mathcal{F}}
\newcommand{\calm}{\mathcal{M}}
\newcommand{\calo}{\mathcal{O}}
\newcommand{\calw}{\mathcal{W}}

%%\newcommand*\O{\mathcal{O}}\CL


%%%%%%%%%%%%%%%%%%%%%%TILDE%%%%%%%%%%%%%
\newcommand{\tcalm}{\tilde{\mathcal{M}}}
\newcommand{\tm}{\tilde{{M}}}
\newcommand{\ddt}{\frac{\partial}{\partial t}}
\newcommand{\tX}{\tilde{{X}}}
\newcommand{\tY}{\tilde{{Y}}}
\newcommand{\tZ}{\tilde{{Z}}}
\newcommand{\tcalf}{\tilde{{\mathcal{F}}}}
\newcommand{\tFd}{\tilde{{F}}_{\Delta}}
\newcommand{\tF}{\tilde{{F}}}
\newcommand{\tD}{\tilde{{D}}}
\newcommand{\tth}{\tilde{h}}
\newcommand{\tcalfd}{\tilde{{\mathcal{F}_{\Delta}}}}

\newcommand{\calfd}{{{\mathcal{F}}}_{\Delta}}
\newcommand{\tcald}{\tilde{{\mathcal{D}}}}

\newcommand{\hcalm}{\hat{\mathcal{M}}}
\newcommand*\dk{{\mathcal{D}\kappa}}
\newcommand*\D{{\Delta}}

%%%%%%%%%%%%%%%%%%%MATHbb%%%%%%%%%
\newcommand*\NN{\mathbb{N}}
\newcommand*\TT{\mathbb{T}}
\newcommand*\RR{\mathbb{R}}
\newcommand*\ZZ{\mathbb{Z}}

%%%%%%%%%%%%%%%%%%%MATHFRAK%%%%%%%%%%%%%%%%
\newcommand{\lien}{{\mathfrak{n}}}
\newcommand{\lieg}{{\mathfrak{g}}}
\newcommand{\lieh}{{\mathfrak{h}}}
\newcommand{\liea}{{\mathfrak{a}}}
\newcommand{\tlien}{\tilde{\mathfrak{n}}}
\newcommand{\lies}{{\mathfrak{s}}}
\newcommand{\oo}{{\mathfrak{o}}}







\newcommand{\R}{{\bf R}}
\newcommand{\BP}{{\bf P}}


\DeclareMathOperator{\POop}{PO}
\newcommand{\PO}{{\POop(2,3)}}

\newcommand{\hm}{{\hat{M}}}
\newcommand{\hx}{{\hat{x}}}
\newcommand{\hz}{{\hat{z}}}


\newcommand{\delx}{{\frac{\partial}{\partial x_1}}}
\newcommand{\dely}{{\frac{\partial}{\partial x_2}}}
\newcommand{\delz}{{\frac{\partial}{\partial x_3}}}

\newcommand{\ka}{{\kappa}}


\let\oldtilde\tilde
\renewcommand*\tilde[1]{\mathchoice{\widetilde{#1}}{\widetilde{#1}}{\widetilde{#1}}{\oldtilde{#1}}}
\let\oldhat\hat
\renewcommand*\hat[1]{\mathchoice{\widehat{#1}}{\widehat{#1}}{\widehat{#1}}{\oldhat{#1}}}

\makeatletter
\let\save@mathaccent\mathaccent
\newcommand*\if@single[3]{%
  \setbox0\hbox{${\mathaccent"0362{#1}}^H$}%
  \setbox2\hbox{${\mathaccent"0362{\kern0pt#1}}^H$}%
  \ifdim\ht0=\ht2 #3\else #2\fi
  }
%The bar will be moved to the right by a half of \macc@kerna, which is computed by amsmath:
\newcommand*\rel@kern[1]{\kern#1\dimexpr\macc@kerna}
%If there's a superscript following the bar, then no negative kern may follow the bar;
%an additional {} makes sure that the superscript is high enough in this case:
\newcommand*\widebar[1]{\@ifnextchar^{{\wide@bar{#1}{0}}}{\wide@bar{#1}{1}}}
%Use a separate algorithm for single symbols:
\newcommand*\wide@bar[2]{\if@single{#1}{\wide@bar@{#1}{#2}{1}}{\wide@bar@{#1}{#2}{2}}}
\newcommand*\wide@bar@[3]{%
  \begingroup
  \def\mathaccent##1##2{%
%Enable nesting of accents:
    \let\mathaccent\save@mathaccent
%If there's more than a single symbol, use the first character instead (see below):
    \if#32 \let\macc@nucleus\first@char \fi
%Determine the italic correction:
    \setbox\z@\hbox{$\macc@style{\macc@nucleus}_{}$}%
    \setbox\tw@\hbox{$\macc@style{\macc@nucleus}{}_{}$}%
    \dimen@\wd\tw@
    \advance\dimen@-\wd\z@
%Now \dimen@ is the italic correction of the symbol.
    \divide\dimen@ 3
    \@tempdima\wd\tw@
    \advance\@tempdima-\scriptspace
%Now \@tempdima is the width of the symbol.
    \divide\@tempdima 10
    \advance\dimen@-\@tempdima
%Now \dimen@ = (italic correction / 3) - (Breite / 10)
    \ifdim\dimen@>\z@ \dimen@0pt\fi
%The bar will be shortened in the case \dimen@<0 !
    \rel@kern{0.6}\kern-\dimen@
    \if#31
      \oldoverline{\rel@kern{-0.6}\kern\dimen@\macc@nucleus\rel@kern{0.4}\kern\dimen@}%
      \advance\dimen@0.4\dimexpr\macc@kerna
%Place the combined final kern (-\dimen@) if it is >0 or if a superscript follows:
      \let\final@kern#2%
      \ifdim\dimen@<\z@ \let\final@kern1\fi
      \if\final@kern1 \kern-\dimen@\fi
    \else
      \oldoverline{\rel@kern{-0.6}\kern\dimen@#1}%
    \fi
  }%
  \macc@depth\@ne
  \let\math@bgroup\@empty \let\math@egroup\macc@set@skewchar
  \mathsurround\z@ \frozen@everymath{\mathgroup\macc@group\relax}%
  \macc@set@skewchar\relax
  \let\mathaccentV\macc@nested@a
%The following initialises \macc@kerna and calls \mathaccent:
  \if#31
    \macc@nested@a\relax111{#1}%
  \else
%If the argument consists of more than one symbol, and if the first token is
%a letter, use that letter for the computations:
    \def\gobble@till@marker##1\endmarker{}%
    \futurelet\first@char\gobble@till@marker#1\endmarker
    \ifcat\noexpand\first@char A\else
      \def\first@char{}%
    \fi
    \macc@nested@a\relax111{\first@char}%
  \fi
  \endgroup
}
\makeatother

\let\oldbar\bar
\renewcommand*{\bar}[1]{\mathchoice{\widebar{#1}}{\widebar{#1}}{\oldbar{#1}}{\oldbar{#1}}}
\let\oldoverline\overline
\renewcommand*{\overline}{\bar}

\newenvironment{smallpmatrix}{\left(\begin{smallmatrix}}{\end{smallmatrix}\right)}
\newcommand*{\sfrac}[2]{#1/#2}

\newcommand*{\dd}{\mathrm{d}}
\newcommand*{\dx}{\dd x}
\newcommand*{\dt}{\dd t}
\newcommand*{\du}{\dd u}
\newcommand*{\dv}{\dd v}
\newcommand*{\dw}{\dd w}


%%%%%%%%%%%%%%%%%%%%%%%%%%%%%%%%%%%%%%%%%%%%%%%%
\graphicspath{{./figures/}}

\newcommand*{\mk}{\mkern -1mu}
\newcommand*{\Mk}{\mkern -2mu}
\newcommand*{\mK}{\mkern 1mu}
\newcommand*{\MK}{\mkern 2mu}

\hypersetup{urlcolor=purple, linkcolor=blue, citecolor=red}

\newcommand*{\romanenumi}{\renewcommand*{\theenumi}{\roman{enumi}}}
\newcommand*{\Romanenumi}{\renewcommand*{\theenumi}{\Roman{enumi}}}
\newcommand*{\alphenumi}{\renewcommand*{\theenumi}{\alph{enumi}}}
\newcommand*{\Alphenumi}{\renewcommand*{\theenumi}{\Alph{enumi}}}
%%%%%%%%%%%%%%%%%%%%%%%%%%%%%%%%%%%%%%%%%%%%%%%%%%

